\documentclass[aspectratio=169,12pt]{beamer}
\usepackage{fontspec}
\usepackage[portuguese]{babel}
\usepackage{xcolor,listings,tikz,booktabs,amsmath,amssymb,graphicx,multicol,array,colortbl}
\usetikzlibrary{arrows.meta,positioning,shapes.geometric,calc}

\definecolor{navy}{HTML}{0D1B3E}
\definecolor{gold}{HTML}{F5A623}
\definecolor{qgreen}{HTML}{1E8B4C}
\definecolor{qblue}{HTML}{1A6EAE}
\definecolor{codebg}{HTML}{EEF3F8}

\usetheme{default}
\useinnertheme{circles}
\setbeamertemplate{navigation symbols}{}
\setbeamertemplate{footline}{%
  \hfill{\color{gray!60}\scriptsize\insertframenumber/\inserttotalframenumber}\hspace{.6em}\vspace{.4em}}

\setbeamercolor{frametitle}{bg=navy,fg=white}
\setbeamercolor{structure}{fg=navy}
\setbeamercolor{alerted text}{fg=gold}
\setbeamercolor{item}{fg=gold}
\setbeamercolor{subitem}{fg=navy}
\setbeamercolor{block title}{bg=qblue,fg=white}
\setbeamercolor{block body}{bg=qblue!8,fg=black}
\setbeamercolor{block title alerted}{bg=gold,fg=navy}
\setbeamercolor{block body alerted}{bg=gold!15,fg=navy}
\setbeamercolor{block title example}{bg=qgreen,fg=white}
\setbeamercolor{block body example}{bg=qgreen!10,fg=black}
\setbeamerfont{frametitle}{size=\large,series=\bfseries}
\setbeamerfont{block title}{series=\bfseries}

\setbeamertemplate{title page}{%
  \begin{tikzpicture}[remember picture,overlay]
    \fill[navy](current page.north west)rectangle(current page.south east);
    \fill[gold]([yshift=.7cm]current page.south west)rectangle(current page.south east);
    \node[anchor=south,yshift=.73cm,navy,font=\tiny\bfseries]at(current page.south)
      {INTENSIVÃO QUANT AI \textbullet{} IMPACT UFSCAR};
  \end{tikzpicture}
  \vspace{.65cm}
  \begin{center}
    {\color{gold}\scriptsize\bfseries INTENSIVÃO QUANT AI \textbullet{} IMPACT UFSCAR}\\[.4cm]
    {\color{white}\LARGE\bfseries\inserttitle}\\[.25cm]
    {\color{gold!85}\normalsize\insertsubtitle}\\[.5cm]
    {\color{white!70}\small\insertauthor}\\[.1cm]
    {\color{white!50}\footnotesize\insertdate}
  \end{center}}

\lstdefinestyle{python}{language=Python,
  basicstyle=\ttfamily\scriptsize\color{qblue},
  keywordstyle=\color{navy}\bfseries,
  stringstyle=\color{qgreen!80!black},
  commentstyle=\color{gray!70}\itshape,
  backgroundcolor=\color{codebg},
  frame=l,framerule=2pt,rulecolor=\color{gold},
  xleftmargin=1em,breaklines=true,showstringspaces=false,tabsize=4,
  extendedchars=false,inputencoding=ascii}
\lstset{style=python}

\author{Felipe Maldonado\\{\scriptsize VP Diretoria Quant~\textbullet{}~Impact UFSCAR}}
\date{Intensivão Quant AI 2026}
\title{Aula 09 --- GenAI, Relatório \& Defesa}
\subtitle{LLMs, Geração de Relatório e Simulado de Defesa}

\begin{document}

\begin{frame}[plain]
\titlepage
\end{frame}

\begin{frame}[fragile]{Nesta Aula}
\begin{columns}[T]
  \column{0.5\textwidth}
\begin{block}{Teoria (20 min)}
\begin{itemize}
  \item Onde GenAI (LLMs) se encaixa de forma prática no pipeline quant
  \item Como redigir prompts estruturados para análise técnica do Claude
  \item Os 7 critérios de avaliação da banca do Itaú Asset
  \item A estrutura do relatório final e regras da defesa oral
\end{itemize}
\end{block}
  \column{0.47\textwidth}
\begin{block}{Código (40 min)}
\begin{itemize}
  \item Setup: wrappers para a API do Claude (\texttt{anthropic} SDK)
  \item \texttt{gerar\_draft\_secao()} — automatizar seções via LLM
  \item \texttt{painel\_tear\_sheet()} — figura consolidada de 5 painéis
  \item Geração automática do rascunho em markdown
  \item Simulação da banca de defesa oral com IA
\end{itemize}
\end{block}
\end{columns}
\end{frame}

\begin{frame}[fragile]{LLMs no Pipeline Quantitativo}
\begin{columns}[T]
  \column{0.5\textwidth}
\begin{itemize}
  \item \textbf{Análise de Resultados}: LLM recebe Sharpe, MDD e volatilidade, interpretando regimes de risco e gerando comentários.
  \item \textbf{Redação Estruturada}: Geração rápida do Sumário Executivo, Metodologia (LaTeX) e Limitações técnicas.
  \item \textbf{Revisão Crítica}: Atuar como advogado do diabo, questionando fragilidades no nosso próprio backtest.
  \item \textbf{Preparação da Defesa}: Simular perguntas difíceis da banca da Itaú Asset.
\end{itemize}
  \column{0.47\textwidth}
\begin{block}{Boas Práticas de Prompt}
\begin{itemize}
  \item Use System Prompts claros (defina persona de analista sênior)
  \item Passe dados estruturados (não mande 'analisar minha estratégia', passe a tabela)
  \item Peça Chain of Thought (ex: 'raciocine sobre regimes de mercado antes de concluir')
  \item Configure temperature baixa para dados numéricos precisos
\end{itemize}
\end{block}\vspace{.2em}
\begin{alertblock}{Segurança e API Keys}
Nunca insira chaves no código. Use variáveis de ambiente (\texttt{ANTHROPIC\_API\_KEY}) e suba apenas o build.
\end{alertblock}
\end{columns}
\end{frame}

\begin{frame}[fragile]{Os 7 Critérios da Banca do Itaú}
\begin{columns}[T]
  \column{0.5\textwidth}
\begin{itemize}
  \item \textbf{1. Conceito da Estratégia (20\%)}: Tese econômica clara (underreaction comportamental, Jegadeesh \& Titman).
  \item \textbf{2. Modelagem (20\%)}: Justificativa dos parâmetros, tratamento de outliers e Markowitz caps.
  \item \textbf{3. Uso de GenAI (15\%)}: Claude usado para análise qualitativa, redação e críticas do relatório.
  \item \textbf{4. Backtest (15\%)}: Sem look-ahead bias, weights.shift(1), custos realistas (0.3\%) e walk-forward.
  \item \textbf{5. Análise de Resultados (15\%)}: Sharpe, MDD, Calmar, Alpha/Beta vs IBOVESPA.
  \item \textbf{6. Conclusão e Próximos Passos (10\%)}: Limitações quantificadas de forma madura.
  \item \textbf{7. Apresentação (5\%)}: Identidade e clareza.
\end{itemize}
  \column{0.47\textwidth}
\begin{alertblock}{Regra de Ouro}
O rigor técnico e a honestidade intelectual sobre as limitações do backtest contam muito mais para a banca do Itaú do que o retorno absoluto.
\end{alertblock}
\end{columns}
\end{frame}

\begin{frame}[fragile]{Pipeline Completo — Resumo}

\begin{center}
\begin{tikzpicture}[node distance=.55cm and .3cm,
  box/.style={draw=navy,fill=navy!8,rounded corners=3pt,
              text width=1.55cm,align=center,font=\tiny\bfseries,inner sep=3pt},
  lbl/.style={font=\tiny,gold},
  arr/.style={-Stealth,navy,thick}]
  \node[box](d){Dados\\yfinance};
  \node[box,right=of d](e){EDA\\análise};
  \node[box,right=of e](s1){Sinal v1\\12-1};
  \node[box,right=of s1](s2){Sinal v2\\vol-adj};
  \node[box,right=of s2](p){Portfólio\\Markowitz};
  \node[box,right=of p](b){Backtest\\EW vs MW};
  \node[box,right=of b](v){WF-OOS\\Custos 0.3\%};
  \node[box,right=of v](g){GenAI\\Claude};
  \node[box,right=of g](r){Relatório\\Defesa};

  \draw[arr](d)--(e);\draw[arr](e)--(s1);\draw[arr](s1)--(s2);
  \draw[arr](s2)--(p);\draw[arr](p)--(b);\draw[arr](b)--(v);
  \draw[arr](v)--(g);\draw[arr](g)--(r);

  \node[lbl,below=.3cm of d]{Aula 2};
  \node[lbl,below=.3cm of e]{Aula 3};
  \node[lbl,below=.3cm of s1]{Aula 4};
  \node[lbl,below=.3cm of s2]{Aula 7};
  \node[lbl,below=.3cm of p]{Aula 6};
  \node[lbl,below=.3cm of b]{Aulas 5-6};
  \node[lbl,below=.3cm of v]{Aula 8};
  \node[lbl,below=.3cm of g]{Aula 9};
  \node[lbl,below=.3cm of r]{Aula 9};
\end{tikzpicture}
\end{center}
\vspace{.2em}
\begin{columns}[T]
  \column{0.48\textwidth}
  \begin{block}{Entregáveis da nossa estratégia}
    \texttt{precos\_ibov} $\cdot$ \texttt{sinal\_v2} $\cdot$ \texttt{pesos\_markowitz}\\
    \texttt{retorno\_walkforward\_liquido} $\cdot$ \texttt{tearsheet\_final.png}
  \end{block}
  \column{0.48\textwidth}
  \begin{exampleblock}{Dica para a Defesa}
    Apresente as limitações (como survivorship bias) antes que a banca pergunte. Isso demonstra maturidade acadêmica e prática.
  \end{exampleblock}
\end{columns}
\end{frame}

\begin{frame}[fragile]{Como Defender — Perguntas Difíceis}
\begin{columns}[T]
  \column{0.5\textwidth}
\begin{block}{Janela 12-1 e Parâmetros}

\textit{``Por que 12-1 e não outra janela?''}\\[.3em]
\textbf{Defesa}: Fixamos o parâmetro baseados na literatura de Jegadeesh \& Titman (1993) antes de ver os dados para evitar overfitting de parâmetros.

\end{block}
\vspace{.2em}
\begin{block}{Equal-Weight vs Otimizado}

\textit{``Por que usar Equal-Weight ou Markowitz restrito?''}\\[.3em]
\textbf{Defesa}: Otimização pura é instável devido a erros de estimação. Impor restrições de peso ($w_i \in [0, 0.20]$) nos protege contra pesos extremos.

\end{block}
  \column{0.47\textwidth}
\begin{block}{Regras da Apresentação}
\begin{itemize}
  \item Nunca invente dados: admita o que não sabe
  \item Foque nos gráficos: a imagem da Tear Sheet final conta a história inteira
  \item Cuidado com o turnover: demonstre que a estratégia sobrevive à fricção de transação
\end{itemize}
\end{block}
\end{columns}
\end{frame}

\begin{frame}[fragile]{O que Vamos Construir}
\begin{columns}[T]
  \column{0.52\textwidth}
\begin{block}{Funções a implementar}
\begin{itemize}
  \item \texttt{chamar\_claude(prompt, system)} $\to$ chamada Anthropic
  \item \texttt{painel\_tear\_sheet(ret\_wf, ret\_ibov)} $\to$ Tearsheet PNG
  \item \texttt{gerar\_relatorio\_completo()} $\to$ exportar draft markdown
\end{itemize}
\end{block}
  \column{0.45\textwidth}
\begin{block}{Entradas (parquets)}
\begin{itemize}
  \item \texttt{retorno\_walkforward\_liquido.parquet}
  \item \texttt{retornos\_mensais\_limpo.parquet}
  \item API key da Anthropic
\end{itemize}
\end{block}
\vspace{.4em}
\begin{exampleblock}{Saídas (parquets)}
\begin{itemize}
  \item \texttt{relatorio\_final\_draft.md}
  \item \texttt{tearsheet\_final.png}
\end{itemize}
\end{exampleblock}
\end{columns}
\end{frame}

\begin{frame}[fragile]{Referências — Pipeline Completo}
\begin{multicols}{2}
\tiny
\begin{itemize}
  \item \textbf{Jegadeesh, N. \& Titman, S.} (1993). Returns to buying winners. \textit{Journal of Finance}, 48(1).
  \item \textbf{Moskowitz, T. et al.} (2012). Time series momentum. \textit{JFE}, 104(2), 228--250.
  \item \textbf{DeMiguel, V. et al.} (2009). Optimal versus naive diversification. \textit{RFS}, 22(5).
  \item \textbf{Bailey, D. \& López de Prado, M.} (2014). The deflated Sharpe ratio. \textit{JPM}, 40(5).
  \item \textbf{Daniel, K. \& Moskowitz, T.} (2016). Momentum crashes. \textit{JFE}, 122(2).
  \item \textbf{Markowitz, H.} (1952). Portfolio Selection. \textit{Journal of Finance}, 7(1).
\end{itemize}
\end{multicols}
\end{frame}


\end{document}
